% ---
% id: statistics/binomial-box/rem:guard
% kind: rem
% title: "Benefit and cost of the honest-containment guard"
% summary: >-
%   The guard buys honest containment for every budget and every allocation,
%   and never raises the failure probability. Its cost is not bounded by one
%   count against the raw quantile interval, yet an endpoint it moves stays
%   within 1/n of the honest value.
% prominence: standard
% status: accepted
% uses: [prop:rule-exact, lem:guard, def:interval, def:endpoint-rule]
% mentions: [thm:monotonicity-limits, thm:box-construction, thm:honest-containment]
% assumes: []
% refs: []
% provenance: newly-derived
% lean:
% ---

\begin{remark}\label{rem:guard}
This remark reads the \tref{honest-containment guard} on the exact rule of
Definition~\ref{def:interval}, whose rule endpoints are the quantiles
$\kappa^{\pm}$; Lemma~\ref{lem:guard} states what holds for every rule. It is
written for a single honest value $p_{\hat c}=\honest(\hat c)$, the
degenerate envelope; for a general envelope every statement below holds side
by side, by the side-wise reading of Definition~\ref{def:endpoint-rule}: what
it says of the lower endpoint with $p_{\hat c}=p^{-}_{\hat c}$, and what it
says of the upper endpoint with $p_{\hat c}=p^{+}_{\hat c}$.

\emph{What the guard buys} is unconditional: the whole envelope interval
$[p^{-}_{\hat c},p^{+}_{\hat c}]$ --- for a single law, $\honest(\hat c)$ ---
lies in $[k^{-}(\hat c)/n,\,k^{+}(\hat c)/n]$ for \emph{every} admissible
budget and \emph{every} admissible allocation, including allocations so
lopsided that the raw quantile interval
$[\kappa^{-}(\hat c),\kappa^{+}(\hat c)]$ excludes the honest value.
Theorem~\ref{thm:box-construction}\,(iii) and Theorem~\ref{thm:honest-containment} rest on exactly that,
and without the guard both would carry a side condition on how large
$\ecomp$ may be. By Proposition~\ref{prop:rule-exact}\,(iii) the widening is
always admissible: it can only lower the failure probability.

\emph{What the guard costs} has two different and equally correct answers,
both stated in Lemma~\ref{lem:guard}\,(iii)--(iv) and
Proposition~\ref{prop:rule-exact}\,(vi); the second alone would overstate how
little the guard costs.

\emph{Against the raw quantile interval} the cost is \eqref{eq:guard-cost},
and it is \emph{not} bounded by one count. A concrete admissible input:
$\Chatcon=\{\hat c\}$, $n=10$, $p_{\hat c}=\tfrac12$,
$\varepsilon^{-}(\hat c)=\tfrac9{10}$ and $\varepsilon^{+}(\hat c)=\tfrac1{20}$,
the allocation summing to $\tfrac{19}{20}$ and so admissible for any
$\ecomp\in[\tfrac{19}{20},1)$. There
$\Pr[\Ncount{\hat c}<7]=\tfrac{53}{64}\le\tfrac9{10}$ while
$\Pr[\Ncount{\hat c}<8]=\tfrac{121}{128}>\tfrac9{10}$, so $\kappa^{-}=7$
against $\lfloor np_{\hat c}\rfloor=5$: the guard widens the interval by two
counts. Proposition~\ref{prop:rule-exact}\,(vi) shows the widening is
unbounded in the same direction --- it is at least $1-p_{\hat c}$ in
frequency, uniformly in $n$, once a single tail's allocation reaches
$1-p_{\hat c}^{\,n}$.

\emph{Against the honest value} the quantity bounded is the emitted
endpoint's \emph{distance} from $\honest(\hat c)$, not the enlargement, and
that distance is at most $1/n$, unconditionally: by \eqref{eq:guard-from-p},
on the side where the guard acts the emitted endpoint sits within $1/n$ of
$\honest(\hat c)$. This is the sense in which the guard is cheap, and it is
the sense that survives the comparison with Theorem~\ref{thm:monotonicity-limits}\,(iv): the
$O(n^{-1/2})$ half-width of \eqref{eq:width} is proved for the
\emph{guarded} endpoints, with no side condition, so the guard never breaks
that bound.
\end{remark}
